\chapter{向量}\label{ch:g10:vectors}

向量刻画一个位移：它有方向和长度，与起点无关。
向量能给出几何事实的简洁证明，一旦引入坐标，
这些问题便化为简单的计算。它们将在
\cref{ch:g10:lines}中成为关键工具，之后还将在
\cref{ch:g11:scal}中用于数量积。

\section{平移与向量}

\begin{definition}[向量]\label{def:g10:vectors:vector}
给定两点 $A$、$B$，\emph{向量}\index{向量} $\vect{AB}$
表示从 $A$ 到 $B$ 的位移，它具有
\emph{方向}（直线 $(AB)$ 的方向，连同
沿直线从 $A$ 指向 $B$ 的走向）和
\emph{长度}（距离 $AB$，也记作 $\norm{\vect{AB}}$）。
两个向量当它们表示相同的位移时\emph{相等}：
$\vect{AB} = \vect{CD}$ 表示，对任一点施加两个位移，
所得结果相同。
\end{definition}

\begin{proposition}[相等向量与平行四边形]\label{prop:g10:vectors:parallelogram}
$\vect{AB} = \vect{CD}$ 当且仅当 $ABDC$（顺序不可颠倒！）为
平行四边形（可以是退化的） --- 等价地，$[AD]$ 与
$[BC]$ 有相同的中点。
\end{proposition}
\admitted

\begin{omfigure}
\begin{tikzpicture}[scale=1.0]
  \coordinate (A) at (0,0);
  \coordinate (B) at (2.2,0.8);
  \coordinate (C) at (1.0,-1.2);
  \coordinate (D) at (3.2,-0.4);
  \draw[omExo, thin] (A) -- (C);
  \draw[omExo, thin] (B) -- (D);
  \draw[-Stealth, omDef, thick] (A) -- (B)
    node[midway, above, font=\small] {$\vect{AB}$};
  \draw[-Stealth, omDef, thick] (C) -- (D)
    node[midway, below, font=\small] {$\vect{CD}$};
  \fill (A) circle (1.5pt) node[left, font=\small] {$A$};
  \fill (B) circle (1.5pt) node[right, font=\small] {$B$};
  \fill (C) circle (1.5pt) node[left, font=\small] {$C$};
  \fill (D) circle (1.5pt) node[right, font=\small] {$D$};
\end{tikzpicture}

{\small 相等向量 $\vect{AB} = \vect{CD}$：方向相同，指向相同，
长度相同。四边形 $ABDC$ 是平行四边形。}
\end{omfigure}

\begin{notation}
当端点无关紧要时，向量常用单个字母表示，
例如 $\vec u$ 或 $\vec v$。\emph{零向量} $\vec 0 = \vect{AA}$ 是
什么也不移动的位移。\emph{相反向量}指
与 $\vect{AB}$ 长度相同、方向相反的向量 $\vect{BA}$，
记作 $\vect{BA} = -\vect{AB}$。
\end{notation}

\section{向量的加法}

\begin{definition}[两个向量的和]\label{def:g10:vectors:sum}
\emph{和}\index{向量！和} $\vec u + \vec v$ 是先作 $\vec u$、
再作 $\vec v$ 所得的位移。
\end{definition}

\begin{theorem}[夏勒关系]\label{thm:g10:vectors:chasles}
对任意三点 $A$、$B$、$C$，有
\[
\vect{AB} + \vect{BC} = \vect{AC}.
\]
\end{theorem}

\begin{proof}
从 $A$ 到 $B$，再从 $B$ 到 $C$ 的位移，把
$A$ 移到 $C$，即 $\vect{AC}$ 所表示的位移。
\end{proof}

\begin{remark}[平行四边形法则]
要将从同一点出发的两个向量 $\vect{AB} + \vect{AC}$ 相加，
补全平行四边形 $ABDC$：其和即为对角线 $\vect{AD}$。
事实上 $\vect{AC} = \vect{BD}$，故
$\vect{AB} + \vect{AC} = \vect{AB} + \vect{BD} = \vect{AD}$，由夏勒
关系可得。
\end{remark}

\begin{omfigure}
\begin{tikzpicture}[scale=1.05]
  \coordinate (A) at (0,0);
  \coordinate (B) at (2.4,0.5);
  \coordinate (C) at (3.4,2.0);
  \draw[-Stealth, omDef, thick] (A) -- (B)
    node[midway, below, font=\small] {$\vec u$};
  \draw[-Stealth, omThm, thick] (B) -- (C)
    node[midway, right, font=\small] {$\vec v$};
  \draw[-Stealth, omMeth, thick] (A) -- (C)
    node[midway, above left, font=\small] {$\vec u + \vec v$};
\end{tikzpicture}\hspace{1.4cm}%
\begin{tikzpicture}[scale=1.05]
  \coordinate (A) at (0,0);
  \coordinate (B) at (2.4,0.5);
  \coordinate (C) at (1.0,1.5);
  \coordinate (D) at (3.4,2.0);
  \draw[omExo, thin] (B) -- (D) -- (C);
  \draw[-Stealth, omDef, thick] (A) -- (B)
    node[midway, below, font=\small] {$\vect{AB}$};
  \draw[-Stealth, omThm, thick] (A) -- (C)
    node[midway, left, font=\small] {$\vect{AC}$};
  \draw[-Stealth, omMeth, thick] (A) -- (D)
    node[pos=0.75, below right, font=\small] {$\vect{AD}$};
  \fill (D) circle (1.2pt) node[right, font=\small] {$D$};
\end{tikzpicture}

{\small 和的两种图示：首尾相接（沙勒关系，左）与
平行四边形法则（右）。}
\end{omfigure}

\begin{example}[利用沙勒关系化简]\label{ex:g10:vectors:chasles}
化简 $\vect{MN} + \vect{NP} + \vect{PQ}$：依次连接各段位移
得 $\vect{MQ}$。化简 $\vect{AB} - \vect{AC}$：将
差改写为与相反向量的和，
\[
\vect{AB} - \vect{AC} = \vect{AB} + \vect{CA}
= \vect{CA} + \vect{AB} = \vect{CB}.
\]
\end{example}

\section{数乘向量}

\begin{definition}[数乘]\label{def:g10:vectors:scalar}
设 $\vec u$ 为非零向量，$k$ 为实数。向量
$k\vec u$\index{向量！数乘}的方向与
$\vec u$ 相同，长度为 $\abs{k} \times \norm{\vec u}$，
并与 $\vec u$ 同向（当 $k > 0$ 时）、反向（当 $k < 0$ 时）。当 $k = 0$ 时，
$0\vec u = \vec 0$。
\end{definition}

\begin{definition}[共线]\label{def:g10:vectors:collinear}
两个向量 $\vec u$ 与 $\vec v$ 称为
\emph{共线}\index{共线向量}，当它们的方向相同时，
即 $\vec v = k \vec u$，其中 $k$ 为某个实数（或其中一个为
$\vec 0$）。
\end{definition}

\begin{remark}
共线是刻画平行与三点共线的向量语言：
\begin{itemize}
  \item 直线 $(AB)$ 与 $(CD)$ 平行，当且仅当 $\vect{AB}$
        与 $\vect{CD}$ 共线；
  \item 三点 $A$、$B$、$C$ 共线，当且仅当 $\vect{AB}$ 与
        $\vect{AC}$ 共线。
\end{itemize}
\end{remark}

\section{向量的坐标}

\begin{definition}[向量的坐标]\label{def:g10:vectors:coords}
在坐标系中，向量 $\vect{AB}$ 的\emph{坐标}
是描述位移的数：
\[
\vect{AB}\,(x_B - x_A,\ y_B - y_A).
\]
向量 $\vec u\,(a, b)$ 使每个点水平移动 $a$ 个单位，并
竖直移动 $b$ 个单位。
\end{definition}

\begin{proposition}[利用坐标进行计算]\label{prop:g10:vectors:coordrules}
设 $\vec u\,(a, b)$ 与 $\vec v\,(c, d)$，且 $k \in \R$。
\begin{enumerate}
  \item $\vec u = \vec v$ 当且仅当 $a = c$ 且 $b = d$；
  \item $\vec u + \vec v$ 的坐标为 $(a + c,\ b + d)$；
  \item $k\vec u$ 的坐标为 $(ka,\ kb)$；
  \item $\norm{\vec u} = \sqrt{a^2 + b^2}$（在标准正交坐标系中）。
\end{enumerate}
\end{proposition}

\begin{proof}
1--3 逐个坐标地说明位移的作用：例如
先作 $\vec u$ 再作 $\vec v$，则点水平移动
先 $a$ 后 $c$，共移动 $a + c$。第 4 条是
将\cref{thm:g10:coordgeom:distance}应用于一条表示
$\vec u$ 的线段所得的距离公式。
\end{proof}

\begin{theorem}[共线判定准则]\label{thm:g10:vectors:det}
两个向量 $\vec u\,(a, b)$ 与 $\vec v\,(c, d)$ 共线
当且仅当
\[
ad - bc = 0 .
\]
\end{theorem}

\begin{proof}
若 $\vec v = k\vec u$，则 $c = ka$，$d = kb$，于是
$ad - bc = a(kb) - b(ka) = 0$。同样地，若 $\vec u = k \vec v$，或
两者之一为零向量，结论同样成立。

反之，设 $ad - bc = 0$ 且 $\vec u \neq \vec 0$，不妨设
$a \neq 0$（$b \neq 0$ 的情形类似）。取 $k = \frac{c}{a}$，则
$c = ka$，由 $ad = bc = b(ka)$，两边除以 $a \neq 0$ 得
$d = kb$，故 $\vec v = k\vec u$。
\end{proof}

\begin{example}\label{ex:g10:vectors:det}
判断 $\vec u\,(3, -2)$ 与 $\vec v\,(-6, 4)$ 是否共线？计算
$ad - bc = 3 \times 4 - (-2)\times(-6) = 12 - 12 = 0$：是，且
$\vec v = -2\vec u$。对 $\vec u\,(3, -2)$ 与 $\vec w\,(1, 5)$，
$3 \times 5 - (-2) \times 1 = 17 \neq 0$：不共线。
\end{example}

\begin{method}[三点共线与平行的证明]\label{met:g10:vectors:alignment}
要证明三点 $A$、$B$、$C$ 共线：
\begin{enumerate}
  \item 计算 $\vect{AB}$ 与 $\vect{AC}$ 的坐标；
  \item 验证\cref{thm:g10:vectors:det}中的判定准则 $ad - bc = 0$；
  \item 得出结论：两向量共线且有公共点 $A$，因此
        三点共线。
\end{enumerate}
用 $\vect{AB}$ 与 $\vect{CD}$ 做同样的计算，可证
直线 $(AB)$ 与 $(CD)$ 平行。
\end{method}

\begin{example}\label{ex:g10:vectors:alignment}
设 $A(1, 2)$、$B(3, 5)$、$C(7, 11)$。则 $\vect{AB}\,(2, 3)$ 与
$\vect{AC}\,(6, 9)$，且 $2 \times 9 - 3 \times 6 = 18 - 18 = 0$，故
三点共线（事实上 $\vect{AC} = 3\vect{AB}$）。
\end{example}

\begin{omfigure}
\begin{tikzpicture}[scale=0.62]
  \draw[thin, gray!40] (-0.4,-0.4) grid (8.4,5.4);
  \draw[-Stealth] (-0.6,0) -- (8.6,0);
  \draw[-Stealth] (0,-0.6) -- (0,5.6);
  \coordinate (A) at (1,1);
  \coordinate (B) at (3,2);
  \coordinate (C) at (7,4);
  \draw[-Stealth, omDef, very thick] (A) -- (B)
    node[midway, above left=-2pt, font=\small] {$\vect{AB}$};
  \draw[-Stealth, omThm, thick] (A) -- (C)
    node[pos=0.8, below right=-2pt, font=\small] {$\vect{AC} = 3\vect{AB}$};
  \fill (A) circle (2pt) node[above left, font=\small] {$A$};
  \fill (B) circle (2pt) node[above left, font=\small] {$B$};
  \fill (C) circle (2pt) node[above left, font=\small] {$C$};
\end{tikzpicture}

{\small 共线点：$\vect{AC}$ 是 $\vect{AB}$ 的数乘。}
\end{omfigure}

\section{练习}

\begin{exercise}[$\star$]\label{exo:g10:vectors:1}
利用 沙勒关系化简：
\[
\vect{AB} + \vect{BD}, \qquad
\vect{KL} + \vect{LM} + \vect{MK}, \qquad
\vect{AC} - \vect{BC}, \qquad
\vect{AB} + \vect{CA}.
\]
\end{exercise}

\begin{exercise}[$\star$]\label{exo:g10:vectors:2}
设 $A(2, 1)$、$B(5, 3)$、$C(-1, 4)$。求
$\vect{AB}$、$\vect{BC}$、$\vect{AC}$ 的坐标，并用坐标检验
$\vect{AB} + \vect{BC} = \vect{AC}$ 成立。
\end{exercise}

\begin{exercise}[$\star$]\label{exo:g10:vectors:3}
设 $\vec u\,(2, -3)$、$\vec v\,(-1, 4)$。求
$\vec u + \vec v$、$3\vec u$、$2\vec u - \vec v$ 的坐标，并计算
$\norm{\vec u}$。
\end{exercise}

\begin{exercise}[$\star$]\label{exo:g10:vectors:4}
分别判断 $\vec u$ 与 $\vec v$ 是否共线：
\[
\text{(a) } \vec u\,(4, 6),\ \vec v\,(6, 9);
\qquad
\text{(b) } \vec u\,(2, -5),\ \vec v\,(-4, 10);
\qquad
\text{(c) } \vec u\,(3, 1),\ \vec v\,(1, 3).
\]
\end{exercise}

\begin{exercise}[$\star$]\label{exo:g10:vectors:5}
设 $A(0, 2)$、$B(4, 0)$、$C(6, 3)$。求点
$D$ 的坐标，使 $ABCD$ 为平行四边形，即
$\vect{AB} = \vect{DC}$。
\end{exercise}

\begin{exercise}[$\star\star$]\label{exo:g10:vectors:6}
设 $A(-2, 1)$、$B(1, 3)$、$C(4, -1)$。
\begin{enumerate}
  \item 求点 $M$ 的坐标，使
        $\vect{AM} = \vect{AB} + \vect{AC}$.
  \item 四边形 $ABMC$ 是什么图形？请说明理由。
\end{enumerate}
\end{exercise}

\begin{exercise}[$\star\star$]\label{exo:g10:vectors:7}
设 $A(1, -1)$、$B(4, 1)$、$C(10, 5)$。判断点 $A$、$B$、$C$
是否共线？对 $A(0, 3)$、$B(2, 2)$、$C(8, -1)$ 作同样的判断。
\end{exercise}

\begin{exercise}[$\star\star$]\label{exo:g10:vectors:8}
设 $A(1, 2)$、$B(5, 4)$、$C(6, 1)$、$D(2, -1)$。证明 $(AB)$ 与
$(CD)$ 平行，再证明 $ABCD$ 为平行四边形。哪一个验证
可推出另一个？
\end{exercise}

\begin{exercise}[$\star\star$]\label{exo:g10:vectors:9}
设 $ABC$ 为三角形，点 $I$ 由
$\vect{AI} = \frac23 \vect{AB}$ 定义。取 $A(0,0)$、$B(6, 3)$、$C(2, 5)$
（使计算简便）：
\begin{enumerate}
  \item 求 $I$ 的坐标；
  \item 设点 $J$ 满足 $\vect{CJ} = \frac23\vect{CB}$，求
        $J$ 的坐标；
  \item 证明 $(IJ)$ 与 $(AC)$ 平行。
\end{enumerate}
\end{exercise}

\begin{exercise}[$\star\star\star$]\label{exo:g10:vectors:10}
设 $ABCD$ 为任意四边形，$I$、$J$、$K$、$L$ 分别为
线段 $[AB]$、$[BC]$、$[CD]$、$[DA]$ 的中点。用坐标
$A(x_A, y_A)$ 等表示，证明 $\vect{IJ} = \vect{LK}$，并由此得出
\emph{任意}四边形各边的中点构成
平行四边形。
\end{exercise}

\section{问题：有向线段的代数}

\begin{problem}[{周末问题 --- 沙勒 体操，
两行重证重心，以及平衡力：
向量究竟有什么用}]
\label{pb:g10:vectors:1}
八年级已证明，三角形的三条中线交于各自
长度的三分之二处 --- 其中藏着一个巧妙的平行四边形
（见初中分册中的重心周末问题）。向量只需
两行即可重证，三线共点随之免费得到。这正是
这门新代数的用途：定理化为有向线段上的计算。
本题训练运算（首先是 沙勒关系），给出两行证明，
并推广瓦里尼翁定理
（\cref{exo:g10:vectors:10}），最后回到向量诞生之处：
力与速度。

\medskip
\noindent\textbf{第一部分 --- 沙勒关系训练。}
\begin{enumerate}
  \item 利用 沙勒关系
        （\cref{thm:g10:vectors:chasles}）化简：
        $\vect{AB} + \vect{BC} + \vect{CD}$； \quad
        $\vect{MA} - \vect{MB}$； \quad
        $\vect{AB} + \vect{CD} + \vect{BC} + \vect{DA}$。
  \item 原点技巧：证明对\emph{任意}点 $O$，
        $\vect{AB} = \vect{OB} - \vect{OA}$。
  \item 设 $I$ 为
        $[AB]$ 的中点，证明以下两种刻画：
        \[
        \vect{IA} + \vect{IB} = \vec 0
        \qquad\text{ 且 }\qquad
        \vect{OI} = \tfrac12\left(\vect{OA} + \vect{OB}\right)
        \ \text{ 对任意 } O .
        \]
  \item 利用 $\vect{AB} = \vect{DC}$，重新求平行四边形的第四个顶点
        $D$，其中 $ABCD$ 为平行四边形，$A(-1, 2)$，
        $B(3, 4)$，$C(6, 0)$。
        并与\cref{pb:g10:coordgeom:1}中用中点技巧所得答案作比较。
  \item 把固定向量 $\vec u$ 加到平面的每一点，
        得到的是哪一种变换 ---
        即初中分册双镜周末问题中用两面镜子发现的那种变换，
        以及初中分册半转周末问题中的两次半转？求
        点 $(1, 2)$ 在 $\vec u = (3, -1)$ 下的像。
\end{enumerate}

\medskip
\noindent\textbf{第二部分 --- 重心，第三种证明。}
定义 $G$ 为三角形 $ABC$ 中满足下列
优美条件的点：
\[
\vect{GA} + \vect{GB} + \vect{GC} = \vec 0 .
\]
\begin{enumerate}[resume]
  \item 利用原点技巧，证明该条件
        唯一确定 $G$：
        $\vect{OG} = \frac13\left(\vect{OA} + \vect{OB} +
        \vect{OC}\right)$ 对任意 $O$ 成立，于是
        $G$ 的坐标是各顶点坐标的\emph{平均数}
        （初中分册重心周末问题中已以数值方式看到这一点）。
  \item 证明 $\vect{AG} = \frac13\left(\vect{AB} +
        \vect{AC}\right)$，并记 $K$ 为
        $[BC]$ 的中点，则 $\vect{AK} = \frac12\left(\vect{AB} +
        \vect{AC}\right)$。由此可得：
        $\vect{AG} = \frac23\,\vect{AK}$ --- 三分之二
        定理的重新证明。
  \item 为何由 $B$ 与 $C$ 引出的中线
        无需再作计算即经过\emph{同一个} $G$？
        比较本定理现有的三种证明 ---
        八年级的隐藏平行四边形、坐标法、向量法 ---
        各用一句话比较。
  \item 数值检验：$A(1, 1)$，$B(5, 2)$，$C(3, 6)$。
        求 $G$，再求 $K$，然后
        验证 $\vect{AG} = \frac23 \vect{AK}$，
        利用共线判别准则（\cref{thm:g10:vectors:det}）。
  \item 物理学中，$G$ 可视为
        顶点处三个等质量质点的平衡点。将 $A$ 处质量加倍（质量
        为 $2, 1, 1$），平衡点变为
        $\vect{OG'} = \frac{2\vect{OA} + \vect{OB} +
        \vect{OC}}{4}$。求第 9 题中三角形的 $G'$，
        并将其位置与 $G$ 相比较。
\end{enumerate}

\medskip
\noindent\textbf{第三部分 --- 共线的应用。}
\begin{enumerate}[resume]
  \item 判断 $\vec u(3, -2)$ 与 $\vec v(-6, 4)$ 是否共线？又
        $\vec w(2, 5)$ 与 $\vec z(4, 9)$ 呢？用
        行列式判别准则判断。
  \item 判断点 $A(1, 2)$、$B(3, 5)$、$C(7, 11)$
        是否共线？
  \item 超越瓦里尼翁定理（\cref{exo:g10:vectors:10}）。在任意
        四边形 $ABCD$ 中，\emph{双中线}连接
        对边的中点（$[IK]$ 与 $[JL]$）。利用
        位置向量（$\vect{OI} = \frac12(\vect{OA} +
        \vect{OB})$ 等），证明两条双中线有
        相同的中点，其位置向量为
        $\frac14\left(\vect{OA} + \vect{OB} + \vect{OC} +
        \vect{OD}\right)$：即四边形的重心。
  \item 用一行推出泰勒斯定理：点 $M$ 与 $N$ 满足
        $\vect{AM} = \frac13\vect{AB}$ 及
        $\vect{AN} = \frac13\vect{AC}$，求
        $\vect{MN}$ 用 $\vect{BC}$ 表示，并得出
        平行关系与比值。
  \item 将第 14 题推广到任意比值 $k$，
        用一句话说明向量演算如何同时包含
        中位线定理
        与泰勒斯定理。
\end{enumerate}

\medskip
\noindent\textbf{第四部分 --- 力与速度。}
\begin{enumerate}[resume]
  \item 一条船受正东方向、速度为 $3$ km/h 的水流推动，
        同时以 $4$ km/h 的速度向正北划行。求合
        速度向量、其模，并描述实际
        航线。（一组非常古老的勾股数组出现了。）
  \item 三个力作用于一个圆环：$\vec u(2, 1)$，
        $\vec v(-3, 2)$，$\vec w(1, -3)$。求合力。
        圆环将如何运动？又为何
        沿任一\emph{闭合}多边形的向量之和总为 $\vec 0$
        （利用 沙勒关系）？
  \item 两个力 $\vec u(4, -1)$ 与 $\vec v(-1, 3)$ 作用于
        一个挂钩上。问需要怎样的第三个力，才能使挂钩保持平衡？
  \item 一名游泳者沿垂直河岸的方向过河，速度为
        $2$ m/s，水流速度为 $1.5$ m/s。求
        合速度及其模。河宽
        $50$ m，问过河需要多长时间，
        游泳者在下游多远处靠岸？
  \item 终章：同一对象，三种身份 --- 位移
        有向线段、坐标对、力或速度。列出
        本题借助向量代数重新证明的
        三个经典定理，并指出本章向量尚无法
        度量的一个几何量 --- 用来
        解决它的工具将随数量积
        在十一年级引入。
\end{enumerate}
\end{problem}
